% define document class and set class options
\documentclass[11pt, a4paper]{article}

%===============================================
% *** LIST OF GENERAL PACKAGES *** %
% language package (french option available as well)
\usepackage[english]{babel}

% input fonts
\usepackage[utf8]{inputenc}
\usepackage[T1]{fontenc}
\usepackage{csquotes}

% customize default page geometry (size of margins, ...)
\usepackage{geometry}
\geometry{hmargin=2cm,vmargin=3cm}
\usepackage{titling}

\usepackage{titlesec}

% writing style
\usepackage{lmodern}
\usepackage{times}	 

% include images (image extensions & paths)
\usepackage{graphicx}
\DeclareGraphicsExtensions{.pdf,.jpg,.png,.eps}
\graphicspath{{figs/}}

\usepackage{wrapfig}

% display hyperlinks in the text
\usepackage{hyperref}  %
\hypersetup{colorlinks,%
            citecolor=red,%
            filecolor=black,%
            linkcolor=blue,%
            urlcolor=blue,%
            breaklinks=true}
            
% customize enumerated lists
\usepackage{enumitem}

% lipsum command to fill in the document (for visualisation purposes)
\usepackage{lipsum}

%===============================================
% *** CAPTIONS AND SUBFIGURES *** %
\usepackage{subcaption}

%===============================================
% *** TABLES *** %
\usepackage{booktabs}
\usepackage{multirow}

%=============================================== 
% *** BIBLIOGRAPHY (FOR BIBLATEX) *** %
\usepackage[
    backend=biber,
    style=ieee,
    ]{biblatex}

\usepackage{bookmark}
\addbibresource{bgtg_lib.bib}

% citation style: apa, ieee, authoryear, ...

%===============================================
% *** MATHS / PHYSICS PACKAGES *** %
\usepackage{amsfonts,amssymb,amsmath,amsthm}
\usepackage{mathrsfs}
\usepackage{mathtools}  % for DeclarePairedDelimiter
\usepackage{siunitx}    % notation physical units

% notations (shortcuts)
\newcommand{\bs}[1]{\boldsymbol{#1}}
\newcommand{\nbpix}{N}
\DeclareMathOperator{\prox}{prox}
\DeclareMathOperator{\sgn}{sgn}
\DeclareMathOperator*{\argmin}{argmin}
\DeclarePairedDelimiter{\norm}{\lVert}{\rVert}

%===============================================
% *** TITLE, AUTHOR AND DATE *** %
% \pretitle{\begin{center}\fontsize{30bp}{30bp}\selectfont}
% \posttitle{\par\end{center}}

% \preauthor{\begin{center}\fontsize{14bp}{14bp}\selectfont}
% \postauthor{\par\end{center}}
%===============================================
% *** MAIN DOCUMENT *** %
\begin{document}

\begin{titlepage}
\centering
% Logos
\begin{minipage}{.25\linewidth}
\includegraphics[width=\linewidth]{../images-figures/ulille.png}
\end{minipage}
\hfill
\begin{minipage}{.25\linewidth}
\centering
\includegraphics[width=\linewidth]{../images-figures/Centrale-Lille-Ecole.png}
\end{minipage}
\hfill
\begin{minipage}{.25\linewidth}
\includegraphics[width=\linewidth]{../images-figures/imt.png}
\end{minipage}

\vfill

% Subtitle

{\large Research project final report \\ \textit{March 2025} \vspace{1\baselineskip}}

% Title
\begin{minipage}{\linewidth}
\huge
\bfseries
\centering
\rule{\linewidth}{1.5pt}\\
Conditional Generation of Bass Guitar Tablature for Guitar Accompaniment in Western Popular Music\\[-3mm]
\rule{\linewidth}{1.5pt}
\end{minipage}

\vfill

% Persons involved 
\begin{minipage}{.45\linewidth}
\textit{Student:}\\
Olivier \textsc{Anoufa}\\
\textit{Data Science Master 2}

\end{minipage}
\hfill
\begin{minipage}{.45\linewidth}
\flushright
\textit{Supervisors:}\\
Alexandre \textsc{D'Hooge}\\
Ken \textsc{Déguernel}\\
\textit{Algomus team, CRIStAL}
\end{minipage}

\vfill

% Abstract
\begin{flushleft}
\justify
\textbf{Abstract ---} The field of symbolic music generation has seen great advancements with the rise of transformer-based architectures.
Addressing a specific need identified through a user study conducted on guitar players, we focus on developing AI tools to generate bass guitar tablatures conditioned on scores of other instruments in Western popular music.
The bass guitar, a crucial component of the rhythmic and harmonic sections in music, presents a unique challenge due to its dual role in providing structure and groove.
Building upon the tablature notation, the most popular way to share music among guitarists, this work adopts modern encoding schemes to integrate tablature representation into transformer-based models.
To this end, the project involves preprocessing a large dataset of music scores and fine-tuning state-of-the-art transformer architectures.
Due to the project being the first of its kind and its creative nature, the evaluation of the generated tablatures is mainly done qualitatively.
\vspace{.5\baselineskip}\\
\end{flushleft}
\vfill


% Period
% October 1st --- March 31st
% \vspace{1\baselineskip}\\
\begin{flushleft}
\textbf{Key-words ---} music information retrieval, conditional generation, guitar accompaniment, bass guitar, tablature;
\vspace{.5\baselineskip}\\
% \textbf{Mots-clés \hspace{1ex}---} effets audio, traitement de l'information musicale, traitement du signal différentiable, reproduction sonore, musique assistée par ordinateur.
\end{flushleft}
\vfill

% More logos
\begin{minipage}{.25\linewidth}
\includegraphics[height=3cm]{../images-figures/logoCRIStAL.png}
\end{minipage}
\hspace{7cm}
\begin{minipage}{.25\linewidth}
\includegraphics[height=3cm]{../images-figures/cnrs.png}
\end{minipage}

\end{titlepage}


% \begin{figure}[t]% We use titling to put a figure on top of the title page
%     \centering
%     \includegraphics[width=.5\linewidth]{logos_empile}
% \end{figure}

% \title{Bass guitar tablature conditional generation for guitar accompaniment in western popular music}

% \author{Olivier \textsc{Anoufa} \\  University of Lille, France \\ Master 2 Data Science: Research project}

% {\let\newpage\relax\maketitle}

\newpage

\section*{Introduction}

% Context and motivation
Natural language processing methods for the generation of symbolic music is a field of research that has seen great development in the last years.
The transformer architecture, introduced by Vaswani and al. in 2017, has been used to generate scores for various instruments, in diverse styles and genres\cite{vaswaniAttentionAllYou2023, leNaturalLanguageProcessing2024}.
An under publication user study led by Bacot and al. showed a potential need for accompaniment generation tools for guitarists.
The questionnaire was answered by 31 guitarists composers, and 7 of them followed up with an interview.
During the interviews, the focus was made on bass guitar lines and drum parts generation without necessarily having to be familiar with the instrument.
Indeed, guitarists generally tend to write basic bass chords to accompany them, and an AI tool could perform this functional task for them\cite{bacot_tablature_2025}.


% Objectives
This need is the starting point of the project that I chose, proposed by the Algomus team which is part of the CRIStAL laboratory.
In this project, we focus on the conditional aspect of symbolic music generation, for an instrument that has not been thoroughly studied yet: the bass guitar.
More precisely, we aim to generate bass guitar tablatures given other instruments' scores, in the context of western popular music.
Our goal will be to try several combinations of instruments and to evaluate the quality of the generated tablatures.


% Abstract challenges: (high-level such as: propose an informatic representation of music...)
This project presents many challenges. At a high level, we first have to scrap and preprocess large amounts of data in the form of music scores.
Then we have to come up with an informatic representation of music that is adapted to the task of generating bass guitar tablatures.
That is, a way to encode music scores in a way that is meaningful to the transformer architecture we will use.



However, adapting the transformer architecture - originally applied on text - to symbolic music presents many challenges.
Tokenization has to be adapted, data is much less available, and attention has to be tailored to the context\cite{leNaturalLanguageProcessing2024}.



% Terms definition (tablature, add the period when it was used)
Tablatures are a way to represent music for string instruments. They contain information about the fingering to use to play the notes and the rhythm.
% add a fig (Kryptonite - 3 Doors down - electric guitar introduction)

% Difference between guitar and bass (role in the music)


% State of the art (survey, DadaGP, Compound Gen, Structure-informed positional encoding, gen coherent drum accompaniment)
% ResearchRabbit to find potential other articles

% Challenges (tokenization, conditional generation, attention, data cleaning), training of the model)

% What we have done (data preprocessing, getting the rythmic instrument, retrieving the tokens, baseline using dadagp (GuitarCTRL))
% draw io fig of the model


% Workplan (what we are going to do)

\newpage

%\nocite{*} % force display of the full content of the .bib file, w/o any citation in the document
\printbibliography% references: print bibliography (with bibtex file)

\end{document}